\emph{Setting and tools.} Throughout, $(M,G)$ satisfies (T1) to (T3) of \claimref{sheaf.tameness}. We use the definitions, tools (a), (b) and the necessity and sufficiency clauses (1), (2) of \claimref{synthesis.wall-rung-count-is-height-of-dying-line-class}: a $k$-rung completion is a smooth symmetric form $h$ on $T_k = T^*M \oplus \underline{\R}^{2k}$ over $\Mp$, nondegenerate on $\Mp$ including $\Sig$, with $h|_{T^*M} = G$ along $\Sig$ and $h = G \perp H_k$ outside the collar $C \cong \Sig \times [-1,1]$; $N(\cdot)$ and $P(\cdot)$ denote maximal negative and positive subbundles. On the physical half of the collar $P(G) \cong P(\overline G) \oplus \shL$ and $N(G) \cong N(\overline G)$, with $\overline G$ the induced form on $T^*M|_\Sig/\shL$ of signature $(p-1, q)$ (\claimref{sheaf.degeneration-sheaf}; proof of \claimref{synthesis.wall-class-one-rung-obstruction}). Line bundles with equal $w_1$ are isomorphic, and $w$ is multiplicative (Fact \claimref{sheaf.imported-real-line-bundles}; \cite[\S 4]{MilnorStasheff1974}). Two classical inputs (Fact \claimref{synthesis.imported-stiefel-connectivity-and-grassmannian-classification}): (S) the Stiefel manifold $V_k(\R^m)$ of $k$-frames is $(m-k-1)$-connected, so a map from a complex of dimension at most $m - k - 1$ into it is null-homotopic; (G) for a complex $X$ of dimension less than $m - k$, homotopy classes of maps $X \to \mathrm{Gr}_k(\R^m)$ correspond to isomorphism classes of rank-$k$ bundles over $X$ \cite[\S 5]{MilnorStasheff1974}. We also use the elementary fact (F) that the space of nondegenerate symmetric forms of signature $(P,Q)$ on $\R^{P+Q}$ retracts onto $\mathrm{Gr}_Q(\R^{P+Q})$ by sending a form to its negative eigenspace, the fibre (forms with a given negative eigenspace) being convex, and that for a fixed form the maximal negative subspaces form a contractible set (graphs of contractions over one of them); so two continuous families of forms of the same signature over a complex are joined by a continuous family of such forms exactly when their negative subbundles are homotopic as sections of the Grassmann bundle, and a continuous join can be smoothed and made to agree with given smooth families near its ends.

\emph{The adjoint frame.} In $\Cl(3,0)$ let $B := e_2 e_3$, so $B\tilde B = 1$ and $B^{2} = -1$: $B$ is the rotor of a half turn about $e_1$, of order four in the rotor group. Left multiplication by $B$ acts freely on the unit rotors $S^{3}$ with quotient $\Sig = L(4,1)$ (writing a rotor as $R = z_1 + z_2\,e_1e_2$ with $z_1, z_2 \in \mathrm{span}\{1, B\}$, it is $(z_1, z_2) \mapsto (Bz_1, Bz_2)$). For a rotor $R$ let $O(R)_{ab} := \langle e_a\, R e_b \tilde R\rangle$; by \claimref{std-algebra.spin-double-cover}, $O(R) \in \mathrm{SO}(3)$ and $O(-R) = O(R)$, and $O(BR) = O(B)O(R) = D\,O(R)$ with $D = \mathrm{diag}(1,-1,-1)$, since $B e_1 \tilde B = e_1$, $B e_{2,3} \tilde B = -e_{2,3}$. A function $f$ on $S^{3}$ with $f(BR) = f(R)$ is a function on $\Sig$, and one with $f(BR) = -f(R)$ is a section of the line bundle $\shL_B$ whose holonomy around the generator of $\pi_1(\Sig) = \mathbb{Z}_4$ is $-1$; $w_1(\shL_B)$ is the nonzero element of $H^{1}(\Sig;\mathbb{Z}_2) = \mathbb{Z}_2$, so $\shL \cong \shL_B$. Hence: the rows $O_1, O_2, O_3$ of $O$ are a global orthonormal frame of the trivial $\underline{\R}^{3}$, with $O_1$ a section of the trivial line and $O_2, O_3$ sections of $\shL$; that is,
\[
\underline{\R} \oplus \shL \oplus \shL \cong \underline{\R}^{3} \quad\text{explicitly, and}\quad \shL \cong \operatorname{span}(O_2) \subset \underline{\R}^{3}. \tag{$*$}
\]
Taking Stiefel-Whitney classes in $(*)$, $(1+x)^{2} = 1$, so $x^{2} = 0$ and $\operatorname{ht}(x) = 2$. Tensoring $(*)$ with $\shL$ (and $\shL \otimes \shL \cong \underline\R$) gives $\shL \oplus \underline{\R}^{2} \cong \shL^{\oplus 3}$. The left-invariant coframe $\sigma_a$ of $S^{3}$ is invariant under left multiplication, so it descends: $T^*\Sig$ is trivial, and since tame walls are two-sided, $T^*M|_\Sig$ and $T_k|_\Sig$ are trivial bundles.

\emph{$\shL \oplus \shL$ is not trivial.} A trivialisation would give a nowhere-zero section, i.e.\ a continuous $f : S^{3} \to \mathrm{span}\{1, B\}\setminus\{0\}$ (the fibre $\R^{2}$ of $\shL \oplus \shL$ read as the plane $\mathrm{span}\{1, B\}$) with $f(BR) = -f(R)$. Since $S^{3}$ is simply connected, $f = |f| e^{B\varphi}$ with $\varphi$ continuous and real; then $\varphi(BR) - \varphi(R) \in \pi + 2\pi\mathbb{Z}$ is a continuous function on a connected space, hence a constant $(2j+1)\pi$, and applying it four times, $0 = \varphi(B^{4}R) - \varphi(R) = 4(2j+1)\pi$, a contradiction.

\emph{(2), the definite side.} If $q = 0$ then $N(\overline G) = 0$ and a two-rung completion would give, by clause (1) of the parent claim, a line bundle $N$ with $\shL \oplus N \cong \underline{\R}^{2}$; then $w_1(N) = x$, so $N \cong \shL$ and $\shL \oplus \shL$ is trivial, which is false. With one rung even less is possible ($x \neq 0$). So $k(\Sig) \geq 3$; and by $(*)$, $\shL$ lies in the trivial bundle of rank three, so clause (2) of the parent claim (the coupling of \claimref{sheaf.one-sided-completion} with twisted collar data, its twist confined to the collar) gives a three-rung completion: $k(\Sig) = 3 = e(\shL)$ ($e(\shL) \geq 3$ by the same argument). Such walls exist: on $M = \Sig \times \R$ with $\ell := O_2 \cdot \sigma$ and the cometric $G_t = A - (1 - t) A(\pi\,\cdot, \pi\,\cdot)$ of the projective-wall construction of the parent claim ($A$ the Euclidean metric in the frame $(\sigma_a, dt)$, $\pi$ the projection on $\ell$), the wall $t = 0$ is tame with dying line $\operatorname{span}(\ell) \cong \shL$, positive definite side $t > 0$. The check builds the three-rung form of clause (2) with $\shL$ embedded as $O_2$ in the null space of $H_3$ and verifies at exact rotors and several $t, \beta$: single valuedness under $R \mapsto BR$, $h|_{T^*M} = G$, $\det h = 2\mu\nu\beta^{2} - t$ and signature $(7,3)$ for $t \geq 0$.

\emph{(1), the indefinite side.} Let $q \geq 1$. Over $\Sig$ choose a complement $V_1$ of $\shL$ in $T^*M|_\Sig$, so $G|_{V_1} = \overline G$, and write the rung summand $\underline{\R}^{4}$ with basis $r_1, \dots, r_4$ (any fixed basis, for instance $a_1, a_2, b_1, b_2$ of $H \perp H$). Over $S^{3}$ put $g_a := \sum_b O_{ab}\, r_{b+1}$ ($a = 1,2,3$): $g_1$ descends to $\Sig$, while $g_2, g_3$ are sections of $\shL \otimes \underline{\R}^{4}$, i.e.\ odd. Let $\theta$ be a local generator of $\shL$; $W := \shL \oplus \underline{\R}^{4}$ has the local frame $(\theta, r_1, g_1, g_2, g_3)$. Define $h_W$ on $W$ by
\[
h_W(\theta,\theta) = 0,\quad h_W(\theta, g_2) = 1,\quad h_W(g_2,g_2) = 0,\quad h_W(g_3,g_3) = -1,\quad h_W(r_1,r_1) = h_W(g_1,g_1) = 1,
\]
all other pairings of frame vectors zero. Under $\theta \mapsto -\theta$ together with $R \mapsto BR$ (which sends $g_2, g_3$ to $-g_2, -g_3$) every entry is unchanged, so $h_W$ is a well-defined form on $W$ over $\Sig$, of signature $(3,2)$: the hyperbolic pair $(\theta, g_2)$ plus $(+,+,-)$. Set $h_0 := \overline G|_{V_1} \perp h_W$ on $T_2|_\Sig = V_1 \oplus W$. Then $h_0|_{T^*M} = \overline G \oplus 0 = G$ along $\Sig$, since $\theta$ spans the radical of $G$ there; $h_0$ is nondegenerate of signature $(p-1,q) + (3,2) = (p+2, q+2)$, the signature of $G \perp H_2$ on the physical side; and
\[
N(h_0) = N(\overline G) \oplus \operatorname{span}(\theta - g_2) \oplus \operatorname{span}(g_3) \cong N(\overline G) \oplus \shL \oplus \shL, \qquad P(h_0) = P(\overline G) \oplus \operatorname{span}(\theta + \tfrac12 g_2) \oplus \operatorname{span}(r_1, g_1) \cong P(\overline G) \oplus \shL \oplus \underline{\R}^{2}.
\]
On the other side, at $\Sig \times \{1\}$ the standard form $h_1 := G \perp H_2$ has $N(h_1) \cong N(\overline G) \oplus \underline{\R}^{2}$ and $P(h_1) \cong P(\overline G) \oplus \shL \oplus \underline{\R}^{2}$, transported to $\Sig$ by tool (b). Here $m := \operatorname{rank} T_2 = 8$ and $\dim \Sig = 3$.

Case $q \geq 2$. Then $P(h_0) \cong P(h_1)$, of rank $p + 2$, and $m - (p+2) = q + 2 \geq 4 > 3$, so by (G) the positive subbundles of $h_0$ and $h_1$ are homotopic as maps $\Sig \to \mathrm{Gr}_{p+2}(\R^{8})$ (in the trivialisation of $T_2|_\Sig$), and by (F), applied to $-h$, $h_0$ and $h_1$ are joined by a family of nondegenerate forms over $\Sig \times [0,1]$.

Case $q = 1$, so $p = 3$ and $N(\overline G)$ is a line bundle, trivial or $\shL$ (these are the only two, as $H^{1}(\Sig;\mathbb{Z}_2) = \mathbb{Z}_2$). If it is trivial, $N(h_0) \cong \underline\R \oplus \shL \oplus \shL \cong \underline{\R}^{3} \cong N(h_1)$ by $(*)$; if it is $\shL$, $N(h_0) \cong \shL^{\oplus 3} \cong \shL \oplus \underline{\R}^{2} \cong N(h_1)$ by $(*)$ tensored with $\shL$. Now $m - 3 = 5 > 3$, so (G) and (F) join $h_0$ to $h_1$ as before. (When $N(\overline G)$ is trivial the frame of $N(h_0)$ is explicit: with $n = (\nu_0, \theta - g_2, g_3)$, $\nu_0$ a unit section of $N(\overline G)$, the vectors $f_a := \sum_b O_{ba}\, n_b$ are invariant under $R \mapsto BR$, $\theta \mapsto -\theta$, so they descend, and the map $\Sig \to V_3(\R^{8})$ they define is null-homotopic by (S); this is the instance the check verifies.)

In both cases let $h(\cdot, t)$, $0 \leq t \leq \frac12$, be the join, reparametrised and smoothed so that it equals $h_0$ at $t = 0$ and agrees with $G_t \perp H_2$ near $t = \frac12$ (that family is nondegenerate for $t > 0$ and is homotopic to $h_1$ through itself), and put $h := G \perp H_2$ for $t \geq \frac12$. This is smooth on $T_2$ over $\Mp$, nondegenerate everywhere including $\Sig$, restricts to $G$ along $\Sig$ and is standard outside $C$: a two-rung completion. With the lower bound of the parent claim, $k(\Sig) = 2 = \operatorname{ht}(x)$.

The check verifies the instance $M = \Sig \times \R$, $G_t = (I - (1-t)\,\ell\ell^{T}) \perp (-dt^{2})$ in the frame $(\sigma_a, dt)$, $\ell = O_2$, which is a tame wall with physical side of signature $(3,1)$ and $N(\overline G) = \operatorname{span}(dt)$: at seven exact rational rotors (not normalised; $O$ is computed as $R e_b \tilde R/(R\tilde R)$ with $GA$), $O$ is orthogonal with $\det O = 1$, $O(BR) = DO(R)$ and $O(-R) = O(R)$; $h_0$ in the fixed basis $(\sigma_1,\sigma_2,\sigma_3,dt,a_1,a_2,b_1,b_2)$ satisfies $h_0(BR) = h_0(R)$, its restriction to $T^*M$ has radical $\ell$ and signature $(2,1,1)$, $h_0$ has signature $(5,3)$ like $G_1 \perp H_2$, the Gram matrix of $(dt, \theta - g_2, g_3)$ is $\mathrm{diag}(-1,-2,-1)$, and the frame $f$ is invariant with Gram of signature $(0,3)$.

\emph{(3), the comparison wall.} On $\Sig' = T^{3}$ with $x$ pulled back from the first circle, the half-angle frame $(c, s) = (\cos\frac\varphi2, \sin\frac\varphi2)$ embeds $\shL$ in $\underline{\R}^{2}$ (the stored torus wall's frame), so clause (2) of the parent claim gives a two-rung completion for the definite cometric $G_t = I - (1-t)\,\ell\ell^{T}$, $\ell = (c, s, 0, 0)$; the check verifies single valuedness under $(c,s) \mapsto (-c,-s)$, $\det h = t - 2\mu\nu$ and signature $(6,2)$ for $t \geq 0$ at rational half-angle points. Both $T^{3}$ and $L(4,1)$ are parallelizable ($w = 1$) and carry $x \neq 0 = x^{2}$, with definite physical sides; the rung counts are $2$ and $3$. So in the definite case the rung count is not a function of $(w(\Sig), x)$, its powers and the signature; by the argument of (2), with a definite side two rungs suffice exactly when $x \neq 0$ and $\shL \oplus \shL$ is trivial.

\emph{Remark (classical, from memory; not used above).} $\shL \oplus \shL$ is the realification of $\shL \otimes \mathbb{C}$, whose first Chern class is the integral Bockstein $\beta x \in H^{2}(\Sig;\mathbb{Z})$, a $2$-torsion class with $x^{2} = \beta x \bmod 2$. On $L(4,1)$, $\beta x = 2 \in \mathbb{Z}_4$; on $T^{3}$, $\beta x = 0$. So the datum that decides the definite case is the torsion class $\beta x$, which the mod-2 class $W(\Sig)$ does not record: a torsion-sensitive invariant in the sense of the wall-class conjecture \claimref{sheaf.wall-class}. Why the invariant has to be a class of the bundle of algebras and not of fibrewise module data: at a point of the wall the Clifford algebra of $G$ is that of a form with a one-dimensional radical, the ideal generated by the dying line is nilpotent (it is spanned by the blades with a factor in the radical, Theorem \claimref{std-algebra.radical-ideal-and-quotient}(i), (ii)), so it acts as zero on every simple module, and the simple modules at the wall are those of the quotient by it, the Clifford algebra of $\overline G$ (Theorem \claimref{std-algebra.radical-ideal-and-quotient}(iii)); no fibrewise module datum sees $\shL$, which enters only through how the fibres are glued. In the language of clause (4) of the parent claim, the jump of the Clifford sheaf across the wall is $\Cl(\shL \otimes \R^{0,2})$, a bundle of graded algebras whose odd part is $\shL \oplus \shL$: it is trivial as a bundle of graded algebras exactly when $\shL\oplus\shL$ is, so with a definite physical side two rungs suffice exactly when the Clifford jump is trivial as a bundle (for $x \neq 0$), while (argued, on $L(4,1)$ and the torus) its graded Morita class, by Donovan and Karoubi (from memory) the clock step with $x^{2}$, is what the indefinite count sees.
